% ============================================================
% Chapter 8 — Level 03: Authorization Testing
% ============================================================
\chapter{Level 03 — Authorization \& Access Control Testing}
\label{ch:level03}

\noteblock{This module is a \textbf{placeholder}. Lab environment: \texttt{https://authz.lab.uzyntra.com} (pending deployment).}

\section{Learning Objectives}
\begin{itemize}
  \item Understand and test for Insecure Direct Object References (IDOR)
  \item Identify Broken Object Level Authorization (BOLA) in APIs
  \item Test for vertical and horizontal privilege escalation
  \item Use Burp Repeater and Intruder for access control testing
\end{itemize}

\section{Vulnerability Explanations}

\subsection{IDOR — Insecure Direct Object Reference}
IDOR occurs when an application uses user-controllable input to access objects directly without verifying that the requesting user is authorized to access that specific object.

\textbf{Classic example:}
\begin{lstlisting}[style=http, caption={IDOR in URL parameter}]
# Authenticated as user 100:
GET /api/orders/1001 HTTP/1.1
Authorization: Bearer <user_100_token>

HTTP/1.1 200 OK
{"orderId": 1001, "userId": 100, "total": "$49.99"}

# Change order ID to another user's order:
GET /api/orders/1002 HTTP/1.1
Authorization: Bearer <user_100_token>

HTTP/1.1 200 OK   <-- IDOR: should be 403
{"orderId": 1002, "userId": 101, "total": "$129.99"}
\end{lstlisting}

\subsection{BOLA — Broken Object Level Authorization}
BOLA is the API-specific name for IDOR, classified as OWASP API Security \#1. The application exposes endpoints that handle object IDs without enforcing authorization per-object.

\subsection{Privilege Escalation}

\begin{table}[H]
\centering
\caption{Types of privilege escalation}
\begin{tabularx}{\textwidth}{>{\bfseries}lX}
\toprule
Type & Description \\
\midrule
Horizontal & Access another user's data at the same privilege level \\
Vertical & Gain higher privilege (user $\to$ admin) \\
\bottomrule
\end{tabularx}
\end{table}

\section{Burp Suite Testing Process}

\subsection{Testing for IDOR}
\begin{enumerate}
  \item Log in as User A — capture a request containing your user/object ID
  \item Send to Repeater
  \item Change the ID to another user's ID (increment by 1, try other values)
  \item If the server returns the other user's data with HTTP 200 $\to$ IDOR confirmed
  \item Log in as User B to confirm the data belongs to User B
\end{enumerate}

\subsection{Automated ID Enumeration with Intruder}
\begin{lstlisting}[style=http, caption={Intruder position for IDOR enumeration}]
GET /api/v1/users/§101§/profile HTTP/1.1
Authorization: Bearer <your_token>
\end{lstlisting}
Set payload: Numbers, from 1 to 500, step 1. Filter by response length — consistent non-zero length = data returned = IDOR.

\section{Module Completion Checklist}

\begin{tcolorbox}[findingstyle, title={Level 03 Burp Checklist}]
\begin{itemize}
  \item[\checkbox] Authenticated as low-privilege user — own object IDs captured
  \item[\checkbox] Object IDs enumerated in Repeater — other user data accessible?
  \item[\checkbox] Intruder used to enumerate full ID range
  \item[\checkbox] Admin endpoints tested with regular user token
  \item[\checkbox] Parameter tampering tested: \texttt{role=admin}, \texttt{isAdmin=true}
  \item[\checkbox] HTTP method tampering tested (GET, POST, PUT, DELETE)
  \item[\checkbox] All findings documented with evidence
\end{itemize}
\end{tcolorbox}
